\section{Primitive quotients on a maximal block system}
Let a transitive string C-group on $km$ letters preserve a maximal
nontrivial block system with $m$ blocks of size $k$. Its block image
$P$ is primitive. Choose an inclusion-minimal distinguished subfamily
$L$ generating $P$, and put $C$ equal to the distinguished generators
outside $L$ commuting with every member of $L$. Denote the remaining
family by $R$.

If the noncommuting diagram of the images of $L$ is disconnected,
its commuting factors are normal in $P$. A primitive nontrivial normal
subgroup is transitive; two commuting transitive subgroups are regular.
Their intersection is central and trivial. Each regular normal factor
must be simple, since a proper characteristic or normal subgroup
would give a block system for the product. Thus $P=T\times T$ in the
left--right action on a nonabelian simple group $T$, and $m=|T|\geq60$.
There are exactly two diagram components. Each irredundant generating
family has length at most $\log_2|T|$, so $|L|\leq2\log_2m$.
There are at most four adjacent positions in $R$.

Otherwise the images of $L$ have a connected diagram. Therefore the
corresponding distinguished vertices occupy an interval and $|R|\leq2$.
Whiston gives $|L|\leq m-1$, with equality only for $P=S_m$.
The primitive block action is not cyclic of odd prime order, since
it is generated by involutions. Except when $m=2$, it consequently
has trivial centralizer. Thus $\gen C$ fixes each block. Its restriction
to one block is faithful: an element trivial on that block and commuting
with $\gen L$ is trivial on all blocks, since $L$ is transitive on them.
Hence $|C|\leq k-1$.

For the disconnected case this gives
\[
 d\leq2\log_2m+k+3<km/2\qquad(m\geq60,\ k\geq2).
\]
For example the difference $km/2-k-3-2\log_2m$ is increasing in $k$
and $m\geq60$, and is positive at $(k,m)=(2,60)$.

Suppose now $k,m>2$ and $L$ is connected. Initially $d\leq k+m$.
Equality would make $P=S_m$ and $\gen C$ a diagonal $S_k$ on the blocks.
The diagram of $C$ is connected because $S_k$ is centerless and directly
indecomposable. Both neighbors of the $L$ interval must be present;
all of $C$ lies on one side of it. The opposite neighbor centralizes
$C$. In the full block stabilizer, the centralizer of this diagonal
$S_k$ is precisely the group $S_m$ permuting the blocks consistently.
Since $L$ lies in this centralizer and maps onto $S_m$, it is the whole
centralizer. The opposite neighbor is therefore already in $\gen L$,
contrary to independence. Consequently
\[
 d\leq k+m-1\leq km/2\quad(k,m\geq3,\ km\geq10).
\]
The last inequality follows from $(k-2)(m-2)\geq2$; the only excluded
pair is $(3,3)$, which is part of the finite degree-nine boundary.

\section{Two-point blocks and the equality marking}
Assume $k=2$ and $m\geq5$. Write $V=\mathbb F_2^m$ for the base
module of $C_2\wr S_m$, $W=G\cap V$, and $P=G/W$.
Any member of $C$ is the constant flip $\lambda=(1,\ldots,1)$,
so $|C|\leq1$. For a connected $L$ we have $|R|\leq2$.
If $P\ne S_m,A_m$, the accepted primitive independent-length bound
$\mu'(P)\leq m-3$ gives $d\leq m$, apart from the primitive degree-six
copy of $S_5$.

The $A_m$-invariant subspaces of $V$, for $m\geq5$, are
$0,\gen\lambda,E,V$, where $E$ is the even-weight space (with the
usual coincident or nested spaces understood). Indeed, from any
nonconstant vector a three-cycle gives a nonzero difference of weight
two. Conjugating by $A_m$, which is transitive on unordered pairs,
generates $E$. A family independent under generation by the group
has at most two members, so the normal-module inequality and
$\mu(A_m)\leq m-2$ give $d\leq m$.

The exceptional degree-six $S_5=PGL_2(5)$ is three-transitive.
Every nonconstant vector of weight one, two or three (or its complement)
has conjugate differences generating $E$; for weight three use
three-transitivity, and for the other cases two-transitivity suffices.
Thus it has the same subspace list. Since $m=6$ is even, every invariant
subspace has module-independent length at most one: an odd vector
generates $V$, a nonconstant even vector generates $E$, and a constant
vector generates its line. The normal-module inequality gives
$d\leq\mu(S_5)+1\leq5<6$.

It remains to consider $P=S_m$. The same module list gives
$\mu_G'(W)\leq2$, so $d\leq m+1$. When $m$ is even the preceding
one-generator argument improves this to $d\leq m$.
Suppose $m$ is odd and $d=m+1$. Equality in the normal-quotient proof
forces $|L|=m-1$. Its block images form a Coxeter path: this is
UB.Extremal for $m\geq7$, and the complete independent-string
certificate for $m=5,6$. They occupy an interval in the full string.
Each of the two outside generators centralizes the block subgroup
omitting at most one endpoint of that path. Its centralizer in $S_m$
is trivial (also for the outer $S_6$ image), so both outside generators
actually belong to $V$.

A vector commuting with every adjacent block transposition except
possibly the first or last is constant off that exceptional endpoint.
The possible nonzero vectors are therefore $\lambda,e_i,\lambda+e_i$,
where $i$ is an endpoint of the block path. The vector $e_i$ generates
all of $V$ under $S_m$; two nonconstant even vectors each generate $E$.
Independence of both outside generators consequently forces them to be
$\lambda$ and $\lambda+e_i$. In particular $W=V$.
The subgroup $\gen L$ has zero intersection with $V$: any nonzero
$S_m$-submodule there contains either $\lambda$ or $E$, making one
of these two outside generators redundant. Thus $\gen L$ is a
complement, the graph of a $1$-cocycle $c:S_m\to V$. The automorphism
\[
 (v,\sigma)\longmapsto(v+c(\sigma),\sigma)
\]
of $V\rtimes S_m$ fixes the base module and sends this complement
to the standard one. As an unordered marked set, up to a group isomorphism carrying the
set of distinguished generators onto the other set, the marking is
exactly
\[
 \lambda,\quad\lambda+e_1,\quad s_1,\ldots,s_{m-1},
 \qquad s_i=(i,i+1).
\]
Its displayed ordered Schlaefli type is $[2,4,3,\ldots,3]$.
The noncentral path may be ordered in either direction, and the isolated
central generator may be placed at either end. These are the four
compatible string orders, forming two classes under global duality,
with representatives $[2,4,3,\ldots,3]$ and $[2,3,\ldots,3,4]$.
We do not identify these two ordered polytopes: their labelled types
differ and neither is the global reverse of the other. They share the
same unordered marked set. In particular the source's page-two type
belongs to the reversed noncentral path, not to its literal displayed
tuple; its page-seven type matches that tuple. The source's uniqueness
claim for the unordered pair $(G,S)$ is retained with this distinction.

For completeness the marking has IP for every odd $m\geq3$.
The central factor $\gen\lambda$ is disjoint from $E\rtimes S_m$.
In the remaining family $v=\lambda+e_1,s_1,\ldots,s_{m-1}$, a subset
of adjacent transpositions generates the Young subgroup $P_T$ on
its interval partition. If $v$ is absent the base part is zero.
If $v$ is present and the initial interval containing $1$ is $A$,
the base part is
\[
 W_A=\operatorname{span}\{\lambda+e_i:i\in A\}.
\]
For proper $A$ the indicated vectors are independent, since their
common outside coordinate determines the sum of their coefficients.
For $A=[m]$ their only relation is the sum of all $m$ vectors, giving
$W_{[m]}=E$. All initial intervals are nested and therefore
$W_A\cap W_B=W_{A\cap B}$. The corresponding Young subgroups satisfy
$P_T\cap P_U=P_{T\cap U}$, because their interval partitions refine
by the intersection of their edge sets. The actual subgroups are
$W_A\rtimes P_T$, so their intersections satisfy IP. Adjoining the
independent central factor preserves it. The full group is transitive
on the $2m$ signed points and is $C_2\wr S_m$.
